\documentclass[11pt]{scrartcl}

\usepackage[stixtwo]{fontsetup}
\usepackage{amsmath}

\begin{document}
\begin{enumerate}
	\item \[
		      a_n
		      = \frac{\left(\sqrt{n^2 + 5n} - n\right)\left(\sqrt{n^2 + 5n} + n\right)}{\sqrt{n^2 + 5n} + n}
		      = \frac{5n}{\sqrt{n^2 + 5n} + n}
		      = \frac{5}{\sqrt{1 + \frac{5}{n}} + 1}
		      \to \frac{5}{2}
	      \]

	\item Minden $n \ge 1$ esetén
	      \(
	      \frac{n}{4} = \frac{n^2}{4n} \le \frac{n^2 + 2}{n + 3} \le \frac{3n^2}{n} = 3n
	      \),
	      így
	      \[
		      1
		      \leftarrow \sqrt[n]{\frac{1}{4}} \sqrt[n]{n}
		      \le b_n
		      \le \sqrt[n]{3} \sqrt[n]{n}
		      \to 1,
	      \]
	      tehát a rendőrelv miatt $b_n \to 1$.

	\item \[
		      c_n
		      = \frac{3^n}{n!} \cdot
		      \frac{1 + \frac{n^2}{3^n}}{1 + \frac{2^n}{n!}}
		      \to 0 \cdot \frac{1 + 0}{1 + 0}
		      = 0,
	      \]
	      mert $\frac{3^n}{n!} \to 0$, $\frac{n^2}{3^n} \to 0$ és
	      $\frac{2^n}{n!} \to 0$.
\end{enumerate}
\end{document}
